\section{Pelaksanaan Proyek Akhir}
\label{sec:praktikum-bab17}

\begin{praktikumbox}
\vspace{1.1em}
\textbf{Judul:} Proyek akhir pemrograman assembly 8086.\\[0.3em]
\textbf{CPMK:} \CPMK{03} dan \CPMK{04} \quad\textbullet\quad
\textbf{Alokasi waktu:} 3 minggu (Minggu 13--15) \quad\textbullet\quad
\textbf{Bentuk:} individu atau berpasangan
\end{praktikumbox}

\subsection{Tujuan}

Menghasilkan satu program assembly yang utuh, benar, teruji, dan
dapat dipertanggungjawabkan sebagai bukti pencapaian CPMK03 dan
CPMK04.

\subsection{Tahapan Pengerjaan}

\begin{enumerate}
    \item \textbf{Minggu 13 --- perencanaan.}
    \begin{enumerate}
        \item Pilih salah satu dari empat pilihan proyek pada
              Bagian~\ref{sec:asesmen-akhir}.
        \item Tulis rancangan satu halaman: tujuan, daftar prosedur,
              masukan/keluaran, dan kasus salah.
        \item Kumpulkan rancangan untuk memperoleh persetujuan.
    \end{enumerate}

    \item \textbf{Minggu 14 --- implementasi dan pengujian.}
    \begin{enumerate}
        \item Implementasikan prosedur satu per satu dan uji segera.
        \item Siapkan minimal dua kasus uji per prosedur: satu benar,
              satu salah.
        \item Catat setiap bug dan cara menemukannya.
        \item Simpan setiap versi \file{.asm} yang berfungsi.
    \end{enumerate}

    \item \textbf{Minggu 15 --- penyelesaian dan presentasi.}
    \begin{enumerate}
        \item Lengkapi komentar kode.
        \item Susun laporan: tujuan, rancangan, kendala, pengujian.
        \item Siapkan demonstrasi maksimal 10 menit.
        \item Tunjukkan satu kasus salah yang ditangani dan jelaskan
              satu keputusan teknis.
    \end{enumerate}
\end{enumerate}

\subsection{Kasus Uji Minimum}

\begin{table}[H]
    \centering
    \caption{Kasus uji minimum yang wajib disiapkan.}
    \label{tab:kasus-uji}
    \small
    \begin{tabular}{@{}p{3.4cm}p{5.2cm}p{4.6cm}@{}}
        \toprule
        \textbf{Jenis kasus} & \textbf{Contoh} &
        \textbf{Perilaku yang diharapkan} \\
        \midrule
        Masukan normal & Nilai dalam rentang wajar &
        Keluaran benar \\
        \addlinespace
        Masukan batas & Nilai terbesar/terkecil &
        Tidak terjadi \eng{overflow} \\
        \addlinespace
        Masukan salah & Huruf pada tempat angka, pembagi nol &
        Ditolak dengan pesan, program tidak berhenti \\
        \addlinespace
        Masukan kosong & Pengguna menekan Enter &
        Tidak membaca sampah penyangga \\
        \bottomrule
    \end{tabular}
\end{table}

\subsection{Hasil yang Dikumpulkan}

\begin{enumerate}
    \item Berkas \file{.asm} proyek yang dapat dirakit tanpa
          kesalahan.
    \item Laporan singkat 2--4 halaman.
    \item Catatan pengujian berisi tabel kasus dan hasil.
    \item Catatan \eng{debugging} berisi bug dan cara penemuan.
    \item Bahan presentasi.
\end{enumerate}

\begin{asesmenbox}
\vspace{1.1em}
\textbf{Asesmen proyek akhir (bobot 15\% dari nilai mata kuliah).}
Penilaian memakai rubrik pada Tabel~\ref{tab:rubrik-proyek}. Contoh
penerapannya ada pada Bagian~\ref{sec:contoh-bab17}.

\begin{tabular}{@{}p{7.4cm}p{5.8cm}@{}}
    \toprule
    \textbf{Kriteria} & \textbf{Bobot terhadap nilai proyek} \\
    \midrule
    Kebenaran fungsi (\CPMK{03}) & 1/5 \\
    \addlinespace
    Kualitas kode (\CPMK{03}) & 1/5 \\
    \addlinespace
    Instruksi dan pengalamatan (\CPMK{02}, \CPMK{03}) & 1/5 \\
    \addlinespace
    Pengujian dan \eng{debugging} (\CPMK{04}) & 1/5 \\
    \addlinespace
    Laporan dan presentasi & 1/5 \\
    \bottomrule
\end{tabular}
\end{asesmenbox}
